\subsection{Degrees}

Consider a ring B and a subring A of B. Endow B with the structure of an (A, A)-bimodule deduced by restriction of scalars from the (B, B)-bimodule structure on \(_{s}B_{d}\) .

PROPOSITION 5. --- Let B be a ring, A a simple subring of B, and S a simple left A-module. Then B is a free left A-module of dimension \(\text{long}_{\text{A}}(\text{B} \otimes_{\text{A}} \text{S})\) .

Let r be the length of A; the A-module \(A_{s}\) is isomorphic to \(S^{r}\) . Now, the left A-module B is isomorphic to \(B \otimes_{A} A_{s}\) (II, §3, No.~4, p.~249), hence to \((\mathrm{B} \otimes_{\mathrm{A}} \mathrm{S})^{r}\) (II, §3, No.~7, p.~255, Proposition 7). We therefore have \(\operatorname{long}_{\mathrm{A}}(\mathrm{B}) = r \operatorname{long}_{\mathrm{A}}(\mathrm{B} \otimes_{\mathrm{A}} \mathrm{S})\) , and Proposition 5 follows from Proposition 3 of VIII, p.~123.

DEFINITION 2. --- Let B be a ring and A a simple subring of B. The dimension of the free left A-module B is called the (left) degree of B over A and is denoted by \(^{(1)}\) \([B:A]_{s}\) .

By replacing A and B with the opposite rings, we deduce from the above that B is a free right A-module. We denote its dimension by \([B : A]_{d}\) and call it the right degree of B over A.

\pandocbounded{\includegraphics[keepaspectratio]{images/part001_1658b52d.jpg}}

We can give an example of a field B and a subfield A such that the degrees \([B:A]_{s}\) and \([B:A]_{d}\) differ. \(^{(2)}\)

Let B be a ring, A a simple subring of B, and S a simple left A-module. Let M be a left A-module and a its length. The A-modules M and \(S^{(a)}\) are isomorphic (VIII, p.~122, Proposition 2), hence so are the B-modules \(B \otimes_{A} M\) and \((\mathrm{B} \otimes_{\mathrm{A}} \mathrm{S})^{(a)}\) . The relation

\[
\mathrm{long} _ {\mathrm{A}} (\mathrm{B} \otimes_ {\mathrm{A}} \mathrm{M}) = [ \mathrm{B}: \mathrm{A} ] _ {s} \mathrm{long} _ {\mathrm{A}} (\mathrm{M})\tag{3}
\]

follows from Proposition 5 and Definition 2.

PROPOSITION 6. --- Let C be a ring, B a simple subring of C, and A a simple subring of B. We then have \([C:A]_{s} = [C:B]_{s}[B:A]_{s}\) .

Let us introduce a basis \((e_{i})_{i\in\mathrm{I}}\) of C viewed as a left B-module and a basis \((f_{j})_{j\in\mathrm{J}}\) of B viewed as a left A-module. Then the family \((f_{j}e_{i})_{j\in\mathrm{J},i\in\mathrm{I}}\) is a basis of C viewed as a left A-module (II, §1, No.~13, p.~222, Proposition 25); Proposition 6 follows.

Remarks. --- 1) Suppose that A is a simple subring of a simple ring B and that the right degree \([B : A]_{d}\) is finite. Let C be the endomorphism ring of B viewed as a right A-module; it is a simple ring by Proposition 4, b) of VIII, p.~123. For any b in B, let \(\gamma(b)\) be the mapping \(x \mapsto bx\) from B to B; then \(\gamma : b \mapsto \gamma(b)\) is an isomorphism from B to a subring of C. Moreover, if \((x_{1}, \ldots, x_{m})\) is a basis of the right A-module B, then the morphism that sends c to \((c(x_{1}), \ldots, c(x_{m}))\) is an isomorphism of left B-modules from C to \(B_{s}^{m}\) ; consequently, we have the relation

\[
[ \mathrm{C}: \gamma (\mathrm{B}) ] _ {s} = [ \mathrm{B}: \mathrm{A} ] _ {d}.\tag{4}
\]

Taking into account relation (2) of VIII, p.~123 applied to the right A-module B, we see that

\[
\mathrm{long} (\mathrm{C}) = [ \mathrm{B}: \mathrm{A} ] _ {d} \mathrm{long} (\mathrm{A}).\tag{5}
\]

\begin{enumerate}
\def\labelenumi{\arabic{enumi})}
\setcounter{enumi}{1}
\tightlist
\item
  Let \(\mathrm{K}\) be a commutative field. If \(\mathrm{A}\) is a simple subalgebra of an algebra \(\mathrm{B}\) of finite degree over \(\mathrm{K}\) , then the left degree of \(\mathrm{B}\) over \(\mathrm{A}\) satisfies the relation \(\left[\mathrm{B}:\mathrm{A}\right]_s[\mathrm{A}:\mathrm{K}] = [\mathrm{B}:\mathrm{K}]\) by Proposition 6 of VIII, p.~125. Likewise, we have \(\left[\mathrm{B}:\mathrm{A}\right]_d[\mathrm{A}:\mathrm{K}] = [\mathrm{B}:\mathrm{K}]\) . The equality \(\left[\mathrm{B}:\mathrm{A}\right]_s = \left[\mathrm{B}:\mathrm{A}\right]_d\) follows.
\end{enumerate}

